\documentclass[12pt]{article}
\usepackage[utf8]{inputenc}
\usepackage[T1]{fontenc}
\usepackage{hyperref}
\title{Y\'el\^i Dnye Language Resources for K-12 Curriculum Development}
\author{Compiled by Claude Sonnet 5.5 (claude-sonnet-5-5) for Andr\'as Kornai}
\date{}
\begin{document}
\maketitle

\section{Introduction}
Y\'el\^i Dnye (ISO 639-3: yle) is the Papuan language of Rossel Island, Papua New Guinea. It is the sole surviving member of its language family. The search found one major scholarly resource, a full reference grammar. It found no school materials.

\section{Primary Educational Resources}
No primers, textbooks or school materials in Y\'el\^i Dnye turned up in the search. The reference description says the language is not known to be taught in schools.

\section{Linguistic and Reference Materials}
Stephen C. Levinson's \emph{A Grammar of Y\'el\^i Dnye: The Papuan Language of Rossel Island} (De Gruyter, 2022, Pacific Linguistics 666) is open access. It covers phonology, morphology, syntax and selected semantic fields. It is written for linguists, in English. Levinson acknowledges the earlier, survey-like work by Henderson (1995). The search found no dictionary.
\begin{itemize}
\item MPI record for the grammar: \url{https://www.mpi.nl/publications/item3372004/grammar-yeli-dnye-papuan-language-rossel-island}
\item OAPEN Library (open-access PDF): \url{https://library.oapen.org/handle/20.500.12657/57212?show=full}
\item De Gruyter open-access listing: \url{https://search.hli.ie/DeGruytereBooksOA/9783110733853}
\end{itemize}

\section{Religious and Community Resources}
No scripture or community-organization resources turned up in the search.

\section{Digital and Audio Resources}
No audio, video, corpus or software resources turned up in the search.

\section{Curriculum Development Context}
The open-access grammar gives a solid base for deriving an orthography description and for training teachers. Pedagogical materials, a dictionary and recordings would still have to be created. The language is spoken by the whole community as a first language, which is a strong basis for mother-tongue schooling.

\end{document}